% p011_d 的电脑重画（等量同种正电荷的电场线；中垂线上 C 点两电荷场强方向（红线）与夹角 θ）
\begin{tikzpicture}[line width=0.7pt, >={Stealth[length=2.2mm]},
  fl/.style={postaction={decorate}, decoration={markings, mark=at position #1 with {\arrow{>}}}}]
  % 电场线：左电荷画一遍，右电荷取镜像
  \foreach \sx in {1,-1} {
    \begin{scope}[xscale=\sx]
      \draw[fl=0.6]  (-1.85,0.3)   .. controls (-1.8,1.0)  and (-1.55,1.5) .. (-1.5,2.2);
      \draw[fl=0.65] (-1.67,0.2)   .. controls (-1.2,0.7)  and (-0.85,1.3) .. (-0.78,2.3);
      \draw[fl=0.7]  (-1.61,0.08)  .. controls (-0.9,0.5)  and (-0.3,1.0)  .. (-0.18,2.3);
      \draw[fl=0.75] (-1.98,-0.29) .. controls (-2.05,-1.1) and (-1.6,-1.7) .. (-1.75,-2.5);
      \draw[fl=0.7]  (-1.67,-0.19) .. controls (-1.2,-0.7) and (-0.8,-1.2) .. (-0.72,-2.25);
      \draw[fl=0.6]  (-1.63,-0.13) .. controls (-1.0,-0.6) and (-0.4,-1.2) .. (-0.1,-1.75);
    \end{scope}
  }
  % 两电荷连线与中垂线（带箭头）
  \draw[fl=0.62] (-1.6,0) -- (0,0);
  \draw[fl=0.62] (1.6,0) -- (0,0);
  \draw[fl=0.43] (0,0) -- (0,2.7);
  \draw[fl=0.75] (0,0) -- (0,-1.7);
  \fill (0,0.7) circle (1.6pt);
  \fill (0,-0.7) circle (1.6pt);
  % 两个正电荷
  \foreach \cx in {-1.9,1.9} {
    \draw[fill=white] (\cx,0) circle (0.3);
    \draw (\cx-0.3,0) -- (\cx+0.3,0);
    \draw (\cx,-0.3) -- (\cx,0.3);
  }
  \node at (-1.1,-0.4) {$A$};
  \node at (-0.22,-0.4) {$B$};
  \node at (-0.42,0.78) {$C$};
  % 红笔：两电荷在 C 点的场强方向及夹角 θ
  \draw[->, red!75!black] (-1.62,0.10) -- (3.0,1.8);
  \draw[->, red!75!black] (1.62,0.10) -- (-3.0,1.8);
  \coordinate (R) at (1.9,0);
  \coordinate (Cc) at (0,0.7);
  \coordinate (O) at (0,0);
  \pic[draw=red!75!black, angle radius=0.9cm, "\red{$\theta$}", angle eccentricity=1.35] {angle = Cc--R--O};
\end{tikzpicture}
